\subsection{Parametric Hand Models}\label{subsec:sota_parametric}

Parametric hand models that encode shape and pose in a compact,
differentiable representation underpin the model-based estimation methods
surveyed in \autoref{subsec:sota_estimation} and constitute the canonical
target representation of the present system.
Among available alternatives, \acrshort{mano}~\cite{MANO:SIGGRAPHASIA:2017}
has achieved broad consensus as the field standard.
Li et al.\ characterise it as the most widely adopted hand model in the
field~\cite{zhang2022nimble}, a standing borne out by its near-universal
integration across the learning-based regressors, analytical
inverse-kinematics pipelines, and motion-capture systems reviewed above.
This adoption rests on a combination of capability and efficiency.
As established in~\autoref{subsec:mano}, the corrective pose blend shapes
substantially reduce the joint-collapse artefacts inherent in plain linear
blend skinning, an advantage that translates directly into metrically
more accurate finger deformation for any downstream application.
Crucially for the present work, Romero et al.\ explicitly note that, because
the blend shapes are linear combinations of vertex offsets, their memory
footprint and runtime are each comparable to the mesh footprint, making
the model suitable for real-time
applications~\cite{MANO:SIGGRAPHASIA:2017}.

NIMBLE~\cite{zhang2022nimble} extends \acrshort{mano}'s surface-only
approach by modelling the hand's internal anatomical structure.
Rather than modelling the hand surface in isolation, NIMBLE assembles the hand
from three anatomical layers: 20 bones modelled as triangular meshes, 7 muscle
groups modelled as tetrahedral meshes, and a single skin mesh, totalling
14{,}970 vertices and 27{,}106 faces~\cite{zhang2022nimble}.
Non-rigid tissue deformation is captured during model construction via a
physically-based registration that incorporates a Neo-Hookean elasticity
term for muscles and skin (bones use rigid deformation only)~\cite{zhang2022nimble},
encoding effects such as muscle bulging and skin wrinkling that lie beyond
the reach of surface-only corrective blend shapes.
NIMBLE targets photorealistic rendering and animation for Metaverse and
virtual reality applications; Li et al.\ demonstrate its visual fidelity
via offline ray-tracing and note integration with rendering engines such
as Blender and Unity~\cite{zhang2022nimble}.
All reported experiments are performed on an
\mbox{NVIDIA GeForce RTX 3090}~\cite{zhang2022nimble}; the system makes no
claim of real-time inference performance.
The substantially larger multi-tissue mesh, the GPU-class hardware required
across all reported experiments, and the design intent of offline photorealistic
rendering place NIMBLE outside the CPU-only, interactive-rate operating point
that constrains this thesis (\autoref{subsec:gap}).

For the present work \acrshort{mano} is therefore adopted as the canonical
parametric target.
Its lightweight forward pass satisfies the real-time CPU constraint directly;
its compact $(\boldsymbol{\theta},\,\boldsymbol{\beta})$ parameterisation
is natively compatible with the analytical inverse-kinematics solver that
forms the core of this system; and its standing as the de facto output
representation across the estimation literature ensures direct comparability
with prior work without additional conversion.
NIMBLE's superior anatomical detail is a genuine advance for applications
where photorealistic mesh fidelity is the priority; under the constraints
of the present setting it offers a level of internal structure that exceeds
what avatar animation requires, at a computational cost that CPU-only
deployment cannot sustain.
